\section{Discussion}
\label{sec:discussion}

\subsection{What drives the results}

Three mechanisms generate everything in Section~\ref{sec:results}. A higher capital exponent raises the marginal product of capital at given factor levels, which raises the interest rate from \genBaseInterestThirtyPct{} per cent in the baseline to \genAnthInterestThirtyPct{} per cent in \genEndYear{}; cohorts with perfect foresight respond by saving more, which raises the capital stock by \genAnthCapitalGapThirtyAbs{} per cent and output with it; and the same higher exponent lowers the labour share one-for-one, which at a unit elasticity pins the wage at $(1-\gamma_t)Y_t/L_t$. Output rises because capital accumulates; wages fall because the share effect dominates the output-per-worker effect. Nothing else is required to account for the paths.

The decomposition also explains why the two scenarios differ as little as they do on wages and as much as they do on output. The automation intensity, and therefore the share effect, is identical in both. All that differs is $Z$, which moves output and so output per worker. The productivity arm's additional \genTfpGainPct{} per cent of technology buys \genAnthOutputPerWorkerGapThirtyAbs{} per cent in output per worker against \genZeroTfpOutputPerWorkerGapThirtyAbs{} per cent, which closes most but not all of the \genLabourShareLogChangePctAbs{} per cent share effect, so wages fall by \genAnthWageGapThirtyAbs{} rather than \genZeroTfpWageGapThirtyAbs{} per cent. On output the same difference compounds through the capital stock and produces GDP gaps of \genAnthGdpGapThirtyAbs{} and \genZeroTfpGdpGapThirtyAbs{} per cent.

\subsection{Reading the fiscal result}

The revenue result --- revenue growing more slowly than output in both arms, and barely at all in the automation-only arm, where it ends \genZeroTfpRevenueGapThirtyAbs{} per cent above baseline against GDP \genZeroTfpGdpGapThirtyAbs{} per cent above --- is the one the OBR flags and does not decompose. In this model it is consistent with composition: labour income falls as a share of output, and the model's effective tax functions imply a higher effective rate on labour income than on capital income. The reported series is total tax revenue, which also carries consumption taxes, and consumption falls \genZeroTfpConsumptionGapThirtyAbs{} per cent in that arm, so the attribution across tax bases is not identified by these runs. That is a feature of the UK tax system rather than of AI, and it would follow from any shock that shifted income from labour to capital by the same amount. What the model adds is the magnitude under these two specific calibrations, and the observation that the sign of the debt effect depends on the spending convention: debt ends \genAnthDebtGapThirtyAbs{} per cent above baseline when government spending is a fixed share of output, and \genFixAnthDebtGapThirtyAbs{} per cent below baseline when spending levels are held --- a reversal that holds in the productivity arm but not in the automation-only arm, where debt still ends \genFixZeroTfpDebtGapThirtyAbs{} per cent above baseline.

The model is not well placed to go further on incidence. It has seven lifetime-earnings types and eighty cohorts, so it can say something about who saves and who works, but it has no household-level benefit rules, no unemployment, and no capital-income distribution beyond cohort assets. The complementary static question --- which households bear a given AI labour-market shock, through which taxes and benefits --- is answered with microsimulation on survey microdata, not here, and the two exercises are deliberately kept separate.

\subsection{What the gamma patch buys}

The patch is not a convenience. With a stepped exponent, the entire level effect of automation arrives in the first period of the transition, so a step run both overstates the \genEndYear{} level gap and understates growth over the window --- the two errors have the same cause, and they are not independent diagnostics. More importantly, a step cannot represent the scenarios that the literature specifies. \citet{korinek2026} build their capability, adoption and task parameters to 2030; \citet{mollimas2026} raise the capital exponent from $1/3$ towards~1 over two decades. Neither experiment was runnable in OG-Core before the patch, and the labour-share path in Section~\ref{sec:results-wages} --- a glide from \genBaseLabourSharePct{} to \genAnthLabourShareThirty{} per cent rather than a jump --- is the direct demonstration that it now is.

Two honest qualifications. The patch does not make the labour-share path informative: at $\varepsilon = \genEpsilon{}$ it remains exactly $1-\gamma_t$, so what the patch delivers is a smooth imposed path rather than a modelled one. And the patch is not upstream in OG-Core, so results produced with it are not reproducible from a stock installation without applying the diff.

\subsection{Where this sits among the published estimates}

The productivity arm's \genAnthGdpGapThirtyAbs{} per cent GDP gain by \genEndYear{} is well below the \genAnthropicTableGdpPct{} per cent of \citet{korinek2026} and above the low end of the range argued for by \citet{acemoglu2025}. That position is a consequence of the model's structure rather than of new evidence: a unit elasticity, no adoption friction, no reinstatement, and a technology gain that arrives on schedule because nothing in the model can delay it. The sensible claim is conditional. If automation raises the UK capital exponent by the amount both scenarios assume, and if measured TFP rises by \genTfpGainPct{} per cent as in the ``substantial'' scenario, then this model says the UK gets \genAnthGdpGapThirtyAbs{} per cent more GDP by \genEndYear{}, almost all of it through capital, with real wages \genAnthWageGapThirtyAbs{} per cent lower and consumption \genAnthConsumptionGapThirtyAbs{} per cent higher. Each ``if'' is an assumption with a named source, and Section~\ref{sec:limitations} lists what the model leaves out that could move any of these numbers in either direction.

\section{Conclusion}
\label{sec:conclusion}

I run two AI scenarios --- automation with a zero technology wedge, motivated by the OBR's March~2026 technological-displacement case, and the ``substantial'' column of \citet{korinek2026} --- through a UK overlapping-generations model with real demographics, eighty cohorts and a calibrated tax system. Both raise the capital exponent from \genGammaBase{} to \genGammaShocked{} over \genStartYear--\genEndYear{}, a \genLabourShareFallPp{} percentage-point fall in the labour share, and differ only in the productivity term solved alongside it. GDP ends \genAnthGdpGapThirtyAbs{} and \genZeroTfpGdpGapThirtyAbs{} per cent above the no-AI path, the capital stock \genAnthCapitalGapThirtyAbs{} and \genZeroTfpCapitalGapThirtyAbs{} per cent above it, the real wage \genAnthWageGapThirtyAbs{} and \genZeroTfpWageGapThirtyAbs{} per cent below it, and consumption \genAnthConsumptionGapThirtyAbs{} per cent above it in the productivity arm but \genZeroTfpConsumptionGapThirtyAbs{} per cent below it in the automation-only arm. Revenue grows more slowly than output in both arms.

Two conclusions survive the model's limitations. The composition of an automation-driven output gain matters as much as its size: in both scenarios the gain is capital accumulation, and neither delivers higher wages or materially higher consumption inside the window. And phasing matters mechanically: representing a phased automation path requires a time-varying capital exponent, which OG-Core did not support, and the patch documented here supplies it in a backward-compatible form that reproduces legacy calibrations exactly. Taking the labour-share mechanism beyond an assumption requires an elasticity of substitution below one and a task structure the single-sector Cobb-Douglas specification does not have; that is the natural next step.
